The \bisect{} redistricting platform's principal defense against gaming
is a cryptographic audit chain linking input data, computation, and output
through SHA-256 hashes.
This paper provides a formal analysis of the audit chain as a gaming defense:
its cryptographic security properties, its three-level provenance architecture,
and its fitness for use as litigation-grade reproducibility evidence under
the Daubert standard.

The SHA-256 collision resistance bound ($2^{80}$ operations under the
birthday paradox) makes undetected manifest manipulation computationally
infeasible.
At $10^{15}$ operations per second on 2026 hardware, generating a
SHA-256 collision requires approximately $3 \times 10^7$ years.

We demonstrate \texttt{bisect label-verify} on three contested states:
North Carolina (2020), Texas (2020), and Wisconsin (2020).
For each state, we show that any party with access to the public manifest
can independently verify the full computation chain from Census Bureau
download through final plan assignment in under 60 seconds on a
standard workstation.

Under Daubert's four criteria (testability, peer review, error rate,
general acceptance), the \bisect{} audit chain satisfies all four:
it is testable by any party, the underlying SHA-256 algorithm is
peer-reviewed and standardized (FIPS PUB 180-4), the error rate
(false non-detection) is bounded by $2^{-80}$, and SHA-256 is
universally accepted in cryptographic practice.

We provide model \dia{} legal language for SHA-256 manifest filing:
the manifest must be filed with the Secretary of State within 30 days
of plan adoption and retained for 10 years.
This language is designed to be directly adoptable by state legislatures
considering algorithmic redistricting statutes.
